\documentclass[10pt,a4paper]{amsart}
\usepackage{amsmath,amssymb,amsthm,graphicx}
\usepackage{verbatim}
\newenvironment{reference}{\comment}{\endcomment}
%\newenvironment{reference}{}{}
\newenvironment{slogan}{\comment}{\endcomment}
\newenvironment{history}{\comment}{\endcomment}

% For commutative diagrams we use Xy-pic
\usepackage[all]{xy}

% We use 2cell for 2-commutative diagrams.
\xyoption{2cell}
% Theorem environments.
%
\theoremstyle{plain}
\newtheorem{theorem}[subsection]{Theorem}
\newtheorem{proposition}[subsection]{Proposition}
\newtheorem{lemma}[subsection]{Lemma}

\theoremstyle{definition}
\newtheorem{definition}[subsection]{Definition}
\newtheorem{example}[subsection]{Example}
\newtheorem{exercise}[subsection]{Exercise}
\newtheorem{situation}[subsection]{Situation}

\theoremstyle{remark}
\newtheorem{remark}[subsection]{Remark}
\newtheorem{remarks}[subsection]{Remarks}

\numberwithin{equation}{subsection}

% Macros
%
\def\lim{\mathop{\mathrm{lim}}\nolimits}
\def\colim{\mathop{\mathrm{colim}}\nolimits}
\def\Spec{\mathop{\mathrm{Spec}}}
\def\Hom{\mathop{\mathrm{Hom}}\nolimits}
\def\Ext{\mathop{\mathrm{Ext}}\nolimits}
\def\SheafHom{\mathop{\mathcal{H}\!\mathit{om}}\nolimits}
\def\SheafExt{\mathop{\mathcal{E}\!\mathit{xt}}\nolimits}
\def\Sch{\mathit{Sch}}
\def\Mor{\mathop{\mathrm{Mor}}\nolimits}
\def\Ob{\mathop{\mathrm{Ob}}\nolimits}
\def\Sh{\mathop{\mathit{Sh}}\nolimits}
\def\NL{\mathop{N\!L}\nolimits}
\def\CH{\mathop{\mathrm{CH}}\nolimits}
\def\proetale{{pro\text{-}\acute{e}tale}}
\def\etale{{\acute{e}tale}}
\def\QCoh{\mathit{QCoh}}
\def\Ker{\mathop{\mathrm{Ker}}}
\def\Im{\mathop{\mathrm{Im}}}
\def\Coker{\mathop{\mathrm{Coker}}}
\def\Coim{\mathop{\mathrm{Coim}}}

% Boxtimes
%
\DeclareMathSymbol{\boxtimes}{\mathbin}{AMSa}{"02}

%
% Macros for moduli stacks/spaces
%
\def\QCohstack{\mathcal{QC}\!\mathit{oh}}
\def\Cohstack{\mathcal{C}\!\mathit{oh}}
\def\Spacesstack{\mathcal{S}\!\mathit{paces}}
\def\Quotfunctor{\mathrm{Quot}}
\def\Hilbfunctor{\mathrm{Hilb}}
\def\Curvesstack{\mathcal{C}\!\mathit{urves}}
\def\Polarizedstack{\mathcal{P}\!\mathit{olarized}}
\def\Complexesstack{\mathcal{C}\!\mathit{omplexes}}
% \Pic is the operator that assigns to X its picard group, usage \Pic(X)
% \Picardstack_{X/B} denotes the Picard stack of X over B
% \Picardfunctor_{X/B} denotes the Picard functor of X over B
\def\Pic{\mathop{\mathrm{Pic}}\nolimits}
\def\Picardstack{\mathcal{P}\!\mathit{ic}}
\def\Picardfunctor{\mathrm{Pic}}
\def\Deformationcategory{\mathcal{D}\!\mathit{ef}}
\usepackage{cleveref}
\crefname{originalreference}{original reference}{original references}
\Crefname{originalreference}{Original reference}{Original references}
\makeatletter
\newlabel{lemma-class-diagonal}{{0FHH}{source}{}{sourcepointer-1}{}}
\newlabel{lemma-class-diagonal@cref}{{[originalreference][0][]0FHH}{[1][0][]source}{}{}{}}
\makeatother
\crefname{originalreference}{original reference}{original references}
\Crefname{originalreference}{Original reference}{Original references}
\makeatletter
\newlabel{lemma-unit}{{0FHF}{source}{}{sourcepointer-2}{}}
\newlabel{lemma-unit@cref}{{[originalreference][0][]0FHF}{[1][0][]source}{}{}{}}
\makeatother
\title{The cohomology of the projective line}
\author{The Stacks Project Authors}
\date{stacks-project:weil.tex:local-snapshot-20261012; GNU FDL 1.2 or later}
\begin{document}
\maketitle
\noindent\textit{Editorial note.} All delivered mathematical text is the author's own, copied from the pinned Stacks Project chapter file weil.tex; no statement or proof step was rewritten, reordered or removed.\par\medskip
\noindent\textit{Editorial note.} The counted proof is the author's complete argument: the computation of the class of the diagonal of the product of the projective line with itself, the bound on the sum of the Betti numbers, and the determination of the dimensions through Poincare duality.\par\medskip
\noindent\textit{Editorial note.} The label lemma-cohomology-P1 names the counted statement, and the proof environment immediately following it in the source is delivered verbatim as its proof.\par\medskip
\noindent\textit{Editorial note.} The counted statement is carried by the author's own lemma environment in the chapter source.\par\medskip
\noindent\textit{Editorial note.} Pre-open gate: python operations/agent-authoring/tools/precard\_gate.py --source weil.tex --label lemma-cohomology-P1 --start 2901 --end 2929 --local-id 28150 returned RESULT=OK (L28150 free, occupancy FREE, public interval clear). Pre-assembly tag self-check: canonical tag resolved and not present in reference/public-items.jsonl.\par\medskip
\noindent\textit{Editorial note.} The statement and proof are copied from the pinned chapter file; the cited inputs are kept as external pointers with their canonical Stacks tags.\par\medskip
\noindent\textit{Allowed standard prerequisites (editorial).}\begin{itemize}
\item The class of the diagonal and the minimal number of pure tensors. Same chapter, the lemma on the pairings of the Kunneth components of the class of the diagonal, implying that the class cannot be written as a sum of fewer pure tensors than the sum of the Betti numbers. Bounds the total dimension of the cohomology of the projective line by two.
\item Nonvanishing of the degree zero cohomology and Poincare duality. Same chapter, the lemma that the degree zero cohomology of a nonempty scheme is nonzero, and axiom (A)(b) on the perfect pairing between complementary degrees. Shows that the dimension in degree zero is nonzero and equals the dimension in the top degree, which pins down the two remaining dimensions.
\item Compatibility of the cycle class map with pullbacks. Axiom (C)(a) on the compatibility of the cycle class map with pullbacks, together with additivity of the cycle class map under addition of cycles. Computes the class of the diagonal of the product of the projective line with itself from the classes of the two rulings.
\end{itemize}\medskip
\ifdefined\hypertarget\hypertarget{sourcepointer-1}{}\fi
\noindent\textit{Source reference.} Original reference 0FHH: sources/stacks-weil/pinned-original-source.tex, source lines 2861--2861. Prerequisite pointer invoked at proof line(s) 2921; it names a standard fact of the same chapter that the counted argument uses, and every step of the counted proof built on it is retained verbatim in the delivered body.\par\medskip
\ifdefined\hypertarget\hypertarget{sourcepointer-2}{}\fi
\noindent\textit{Source reference.} Original reference 0FHF: sources/stacks-weil/pinned-original-source.tex, source lines 2800--2800. Prerequisite pointer invoked at proof line(s) 2926; it names a standard fact of the same chapter that the counted argument uses, and every step of the counted proof built on it is retained verbatim in the delivered body.\par\medskip
\begin{lemma}
\label{lemma-cohomology-P1}
Assume given (D0), (D1), (D2), and (D3) satisfying (A), (B), and (C).
Then $H^*(\mathbf{P}^1_k)$ is $1$-dimensional in dimensions $0$ and $2$
and zero in other degrees.
\end{lemma}

\begin{proof}
Let $x \in \mathbf{P}^1_k$ be a $k$-rational point. Observe that
$\Delta = \text{pr}_1^*x + \text{pr}_2^*x$ as divisors on
$\mathbf{P}^1_k \times \mathbf{P}^1_k$. Using axiom (C)(a)
and additivity of $\gamma$ we see that
$$
\gamma([\Delta]) =
\text{pr}_1^*\gamma([x]) +
\text{pr}_2^*\gamma([x]) =
\gamma([x]) \otimes 1 + 1 \otimes \gamma([x])
$$
in $H^*(\mathbf{P}^1_k \times \mathbf{P}^1_k) =
H^*(\mathbf{P}^1_k) \otimes_F H^*(\mathbf{P}^1_k)$.
However, by Lemma \ref{lemma-class-diagonal}
we know that $\gamma([\Delta])$ cannot be written
as a sum of fewer than $\sum \beta_i$ pure tensors
where $\beta_i = \dim_F H^i(\mathbf{P}^1_k)$.
Thus we see that $\sum \beta_i \leq 2$.
By Lemma \ref{lemma-unit} we have $H^0(\mathbf{P}^1_k) \not = 0$.
By Poincar\'e duality, more precisely axiom (A)(b),
we have $\beta_0 = \beta_2$. Therefore the lemma holds.
\end{proof}

\end{document}
